\section{Comparaison des campagnes (ancienne machine vs campagne locale)}

\subsection{Objectif}

Les sections précédentes ont été mises à jour avec des mesures locales (cette machine).
Cette section compare les indicateurs clés avec les valeurs historiques présentes dans la version
précédente du rapport (campagnes d'une autre machine et d'un état de code antérieur), afin
d'expliquer pourquoi certaines performances augmentent ou diminuent.

\subsection{Comparaison synthétique}

\begin{table}[H]
\centering
\caption{Écarts entre ancienne campagne et campagne locale}
\label{tab:campagne-compare}
\begin{tabular}{@{}lrrr@{}}
\toprule
\textbf{Indicateur} & \textbf{Ancienne campagne} & \textbf{Campagne locale} & \textbf{Écart} \\
\midrule
Baseline \(K0\) (ms/iter) & 3.979834 & 8.041472 & \(+102.1\%\) \\
Vectorisation \(\Delta_{\%}(K0)\) SoA vs AoS & \(+0.034\%\) & \(+6.179\%\) & \(+6.145\) points \\
OpenMP \(S(8)\) & 1.697653 & 2.436151 & \(+43.5\%\) \\
OpenMP \(K0@1\) (ms/iter) & 2.555276 & 8.684056 & \(+239.8\%\) \\
\bottomrule
\end{tabular}
\end{table}

\subsection{Interprétation des écarts}

\begin{enumerate}[leftmargin=*]
  \item \textbf{Différence de machine et d'environnement d'exécution.}
  Les performances absolues dépendent fortement du CPU, de la mémoire, du niveau de charge de la
  machine et de l'environnement (ici WSL2). Il est donc normal que les temps absolus changent
  sensiblement entre deux machines.

  \item \textbf{Évolution du code entre campagnes.}
  Cette version n'est pas strictement identique à l'ancienne:
  \begin{itemize}[leftmargin=*]
    \item le backend OpenMP et le backend serial SoA partagent maintenant un coeur d'avance commun;
    \item \(K2/K3\) sont activés en permanence comme \emph{work-time}, ce qui ajoute de
    l'instrumentation interne et augmente les temps absolus;
    \item le backend OMP conserve un coût structurel à 1 thread (buffers touched, merge, dedup, replay).
  \end{itemize}
  Ces changements améliorent la comparabilité méthodologique, mais ils modifient les valeurs absolues.

  \item \textbf{Différence de protocole selon les sections.}
  Baseline/Vectorisation sont mesurés avec rendu actif (\texttt{measure\_temps.sh}), alors que la
  section OpenMP est mesurée en \texttt{--no-render} (\texttt{measure\_openmp.sh}). Les chiffres
  absolus entre sections ne sont donc pas directement comparables.
\end{enumerate}

\subsection{Ce qui reste stable malgré les écarts}

\begin{itemize}[leftmargin=*]
  \item Le hotspot principal reste \(K1\) (avance des fourmis).
  \item La vectorisation SoA seule ne garantit pas un speedup en mode séquentiel.
  \item OpenMP apporte un speedup net avec le nombre de threads (\(S(8)>S(4)>S(2)>S(1)\)).
\end{itemize}
